%% Question bank chapter file - generated by bank_tools.py build (do not edit while the bank is still being built:
%% it is regenerated from the per-test data in Question Bank\data). Every question is one bankq block:
%%   \begin{bankq}[id=..., source={...}, paper=.., qno=.., type=..., subtopic={...}, answer={...}, answer_src=key|solved,
%%                 working={...}, difficulty=1-10, diff_parts={K S I C F}, diff_note={...},
%%                 group=<passage id or empty>, also_in={duplicate copies in other tests}]
%% difficulty = min(10, K+S+I+C+F-1): K concepts 1-3, S steps 1-4, I insight 0-2, C computation 0-2, F format 0-1
%%   ...question LaTeX, options, figures (\bfig / \bfigi -> img/)...
%%   \end{bankq}
%% Passages shared by several questions are bankpassage blocks; their questions carry group=<passage id>.
\documentclass[10pt]{article}
\usepackage[a4paper,margin=1.8cm]{geometry}
\usepackage[T1]{fontenc}
\usepackage{amsmath,amssymb,mathtools}
\usepackage{graphicx}
\usepackage[export]{adjustbox}
\usepackage{xcolor,needspace,keyval}
\usepackage{enumitem}
\usepackage{booktabs,array,tabularx}
\usepackage[hidelinks]{hyperref}
\setlength{\parindent}{0pt}
\setlength{\parskip}{3pt}
\makeatletter
\define@key{bq}{id}{\def\bq@id{#1}}\define@key{bq}{source}{\def\bq@src{#1}}\define@key{bq}{paper}{\def\bq@paper{#1}}
\define@key{bq}{qno}{\def\bq@qno{#1}}\define@key{bq}{type}{\def\bq@type{#1}}\define@key{bq}{subtopic}{\def\bq@sub{#1}}
\define@key{bq}{answer}{\def\bq@ans{#1}}\define@key{bq}{answer_src}{\def\bq@asrc{#1}}\define@key{bq}{working}{\def\bq@work{#1}}
\define@key{bq}{group}{\def\bq@group{#1}}\define@key{bq}{also_in}{\def\bq@also{#1}}
\define@key{bq}{difficulty}{\def\bq@diff{#1}}\define@key{bq}{diff_parts}{\def\bq@dparts{#1}}\define@key{bq}{diff_note}{\def\bq@dnote{#1}}
\newenvironment{bankq}[1][]{\def\bq@work{}\def\bq@also{}\def\bq@group{}\def\bq@diff{?}\def\bq@dparts{}\def\bq@dnote{}\setkeys{bq}{#1}%
  \par\needspace{4\baselineskip}\medskip\noindent{\small\textbf{\bq@id}\hfill\color{gray}[\bq@src, Paper \bq@paper, Q\bq@qno]
  \bq@type\ $\cdot$ \bq@sub\ $\cdot$ \textbf{D\bq@diff}}\par\nopagebreak\relax}
  {\par{\small\color{teal}Answer: \textbf{\bq@ans}\ifx\bq@work\@empty\else\ (solved: \bq@work)\fi
  \ifx\bq@also\@empty\else\par\color{gray}Also in: \bq@also\fi}\par}
\define@key{bp}{id}{\def\bp@id{#1}}\define@key{bp}{source}{\def\bp@src{#1}}\define@key{bp}{questions}{\def\bp@qs{#1}}
\newenvironment{bankpassage}[1][]{\setkeys{bp}{#1}\par\needspace{4\baselineskip}\bigskip\noindent
  {\small\textbf{Passage \bp@id}\hfill\color{gray}[\bp@src; questions \bp@qs]}\par\nopagebreak\relax\itshape}{\par}
\makeatother
\newcommand{\bfig}[2]{\par\smallskip{\centering\includegraphics[width=#1pt,max width=0.62\linewidth,max height=0.32\textheight]{img/#2}\par}\smallskip}
\newcommand{\bfigi}[2]{\includegraphics[width=#1pt,max width=\linewidth,max height=0.32\textheight,valign=c]{img/#2}}
\newcommand{\bmissing}[1]{\textit{[figure missing in the source conversion: \detokenize{#1}]}}
\newenvironment{qoptions}{\begin{itemize}[leftmargin=1.8em,itemsep=1pt,topsep=2pt,label={}]}{\end{itemize}}
\newcommand{\opt}[2]{\item[(#1)] #2}
\providecommand{\fwstart}{}\providecommand{\fwend}{}\providecommand{\qfullsplit}{}
\newcolumntype{L}[1]{>{\raggedright\arraybackslash}p{#1}}
\newcolumntype{C}[1]{>{\centering\arraybackslash}p{#1}}
\newcolumntype{R}[1]{>{\raggedleft\arraybackslash}p{#1}}
\providecommand{\sectionhdr}[1]{}\providecommand{\subsectionhdr}[1]{}


\begin{document}
{\Large\bfseries Electrochemistry --- Question Bank}\par\smallskip
{\small Chemistry chapter C10. Questions from all converted test papers, grouped by source series. Answers: printed key where it matches the paper, otherwise solved (answer\_src=solved, with working). Exact repeats across papers appear once, with also\_in listing the other copies.}\par\medskip\hrule

\section*{Score-I papers}
\begin{bankpassage}[id=SI12-P2-CHEM-S1-G7-8, source={Score-I - Test 12 (09 Mar 2022)}, questions={7-8}]
$E^0_{\text{MnO}_4^-/\text{Mn}^{2+}}=1.51$ V; \ $E^0_{\text{Zn}^{2+}/\text{Zn}}=-0.76$ V\\
$E^0_{\text{Cl}_2/\text{Cl}^-}=1.36$ V; \ $E^0_{\text{Ag}^+/\text{Ag}}=0.8$ V\\
$E^0_{\text{Fe}^{3+}/\text{Fe}^{2+}}=0.77$\\
$K_{sp}(\text{AgCl})=10^{-10}$\\
$\dfrac{2.303RT}{F}=0.06$, \ ($\log2=0.3$)
\begin{center}\small
\begin{tabular}{|p{0.9cm}|p{8.8cm}|p{0.9cm}|p{1.8cm}|p{0.9cm}|p{1.8cm}|}
\hline
 & \centering Column-I & & \centering Column-II & & \centering Column-III\tabularnewline
\hline
(I) & Pt $|$ H$_{2(g)}$(2atm) $|$ H$^+$(0.1M) $\|$ Cl$^-$(0.1) $|$ Cl$_{2(g)}$ (2atm)$|$ Pt & (i) & 0.6 V & (P) & 0.48\\
(II) & Zn(s) $|$ Zn$^{+2}$(0.01M) $\|$ Ag$^+$(0.1M) $|$ Ag(s) & (ii) & 1.56 & (Q) & 1.56\\
(III) & Pt $|$ Fe$^{+2}$(1M), Fe$^{+3}$(1M) $\|$ MnO$_4^-$(1M), H$^+$(pH =2), Mn$^{+2}$(1M) $|$ Pt & (iii) & 0.74 V & (R) & 0.26\\
(IV) & Ag(s) $|$ AgCl(s) $|$ Cl$^-$(0.1M) $\|$ Ag$^+$ (0.1 M) $|$ Ag(s) & (iv) & 1.36 V & (S) & 1.498\\
\hline
\end{tabular}
\end{center}
\end{bankpassage}
\begin{bankq}[id=SI12-P2-CHEM-S1-Q07, source={Score-I - Test 12 (09 Mar 2022)}, paper=2, qno=7, type=single, subtopic={Nernst equation}, answer={D}, answer_src=key, difficulty=3, diff_parts={K2 S2 I0 C0 F0}, diff_note={concepts: Nernst equation, cell EMF and spontaneity | steps: E of cell I; E of cell II; E of cell III and IV; match the properties}, group=SI12-P2-CHEM-S1-G7-8]
Select the correct option
\begin{qoptions}
\opt{A}{I -- (i) -- R}
\opt{B}{II -- (iii) -- S}
\opt{C}{III -- (II) -- Q}
\opt{D}{IV -- (i) -- P}
\end{qoptions}
\end{bankq}
\begin{bankq}[id=SI12-P2-CHEM-S1-Q08, source={Score-I - Test 12 (09 Mar 2022)}, paper=2, qno=8, type=single, subtopic={Nernst equation}, answer={A}, answer_src=key, difficulty=3, diff_parts={K1 S2 I0 C1 F0}, diff_note={concepts: Nernst equation | steps: E of cell I; E of cell II; E of cell III; E of cell IV | C: four Nernst evaluations}, group=SI12-P2-CHEM-S1-G7-8]
Lowest value of $E_{\text{cell}}$ is for Cell
\begin{qoptions}
\opt{A}{III}
\opt{B}{I}
\opt{C}{II}
\opt{D}{IV}
\end{qoptions}
\end{bankq}

\section*{TOAS papers}
\begin{bankq}[id=TOAS16-P2-CHEM-S1-Q03, source={TOAS - Test 16 (04 Sep 2022)}, paper=2, qno=3, type=multi, subtopic={Nernst equation}, answer={C,D}, answer_src=key, difficulty=5, diff_parts={K2 S2 I0 C1 F1}, diff_note={concepts: Nernst equation, solubility product | steps: concentration of lead ion from Ksp; standard emf; correction term; spontaneity}]
Consider the cell Pb(s) $|$ PbI$_2$(s) $|$ I$^-$(aq. 0.1M) $\|$ Cu$^{2+}$(aq. 0.01M) $|$ Cu(s) at temperature $T$.\par
Given $K_{\text{sp PbI}_2}=10^{-8}$ and $E^\circ_{\text{Pb}^{2+}/\text{Pb}}=-0.13$V, $E^\circ_{\text{Cu}^{2+}/\text{Cu}}=0.34$V and $\dfrac{2.303\,RT}{F}=0.06$V.\par
Select the correct option(s) :
\begin{qoptions}
\opt{A}{Above cell is nonspontaneous}
\opt{B}{$E_{\text{cell}}=-0.59$ V}
\opt{C}{$[\text{Pb}^{2+}]_{\text{Anode}}=10^{-6}$ M}
\opt{D}{$E_{\text{cell}}=0.59$ V.}
\end{qoptions}
\end{bankq}
\begin{bankq}[id=TOAS18-P1-CHEM-S1-Q03, source={TOAS - Test 18 (30 Oct 2022)}, paper=1, qno=3, type=multi, subtopic={Nernst equation}, answer={A,B,C,D}, answer_src=solved, working={Cell reaction $\mathrm{Ag^++Cl^-\to AgCl}$: $E^0>0$, $\Delta G^0<0$, $\Delta S<0$ so $dE/dT<0$; diluting the chloride side lowers $E$}, difficulty=3, diff_parts={K1 S2 I0 C0 F1}, diff_note={concepts: temperature and concentration dependence of cell emf | steps: overall reaction; sign of E0 and delta G; temperature coefficient from delta S; Nernst effect of dilution}]
Select correct statement(s) for the cell with notation :
\par\noindent Ag $|$ AgCl $|$ HCl $||$ AgNO$_3$ $|$ Ag ; $(\Delta S)_{\text{reaction}}=-$ve
\begin{qoptions}
\opt{A}{$E^0_{\text{cell}}>0$}
\opt{B}{$\Delta G^0<0$}
\opt{C}{$E_{\text{cell}}$ decreases by increasing temperature}
\opt{D}{$E_{\text{cell}}$ decreases by diluting anodic solution}
\end{qoptions}
\end{bankq}
\begin{bankq}[id=TOAS18-P1-CHEM-S2-Q09, source={TOAS - Test 18 (30 Oct 2022)}, paper=1, qno=9, type=single, subtopic={Cells}, answer={C}, answer_src=solved, working={P: concentration cell, spontaneous, gives work, anodic side gains ions; Q: both sides 1 N in H+, a concentration-type cell with E = 0; R: spontaneous and does work; S: spontaneous, work, chloride falls at the anode}, difficulty=4, diff_parts={K2 S2 I1 C0 F0}, diff_note={concepts: electrochemical cells, concentration cells | steps: nature of each cell; sign of the emf; spontaneity; concentration changes}]
Answer the following by appropriately matching the lists based on the information given in the list.
\begin{center}\footnotesize\setlength{\tabcolsep}{2pt}
\begin{tabular}{|L{0.6cm}L{3.0cm}|L{0.6cm}L{2.7cm}|}\hline
&List-I&&List-II\\\hline
(P)&Zn$|$\allowbreak Zn$^{2+}$(C)$||$\allowbreak Zn$^{2+}$ (2C)$|$\allowbreak Zn&(1)&Spontaneous cell reaction\\\hline
(Q)&H$_2$(P)$|$\allowbreak HCl(1N)$||$\allowbreak H$_2$SO$_4$(1N)$|$\allowbreak H$_2$(P)&(2)&According to cell representation electric work is produced.\\\hline
(R)&Cu$|$\allowbreak Cu$^{2+}$(0.01M)$||$\allowbreak Ag$^+$(0.1M)$|$\allowbreak Ag&(3)&Concentration of cation in cathodic compartment increases upto equilibrium\\\hline
(S)&Ag$|$\allowbreak AgCl$|$\allowbreak (KCl,0.1M)$||$\allowbreak Ag$^+$(0.01M)$|$\allowbreak Ag&(4)&Concentration cell\\\hline
&&(5)&Concentration of Cl$^-$ in anodic compartment decreases\\\hline
\end{tabular}\end{center}
Correctly match among the following is :

[Given : $E^\circ_{\mathrm{Ag^+/Ag}}=0.8$ V, $E^\circ_{\mathrm{Cu^{+2}/Cu}}=0.34$ V, $K_{sp}$(AgCl) $=10^{-10}$]
\begin{qoptions}
\opt{A}{P $\to$ 1,3;Q $\to$ 3,4;R $\to$ 1,2;S $\to$ 1,2,3,4}
\opt{B}{P $\to$ 1,2,4;Q $\to$ 3,5;R $\to$ 1,3;S $\to$ 1,2,4}
\opt{C}{P $\to$ 1,2,4;Q $\to$ 4;R $\to$ 1,2;S $\to$ 1,2,5}
\opt{D}{P $\to$ 1,2,3;Q $\to$ 3,4;R $\to$ 1,2,5;S $\to$ 1,2,3,5}
\end{qoptions}
\end{bankq}
\begin{bankq}[id=TOAS18-P1-CHEM-S3-Q03, source={TOAS - Test 18 (30 Oct 2022)}, paper=1, qno=3, type=numerical, subtopic={Conductance}, answer={4}, answer_src=solved, working={$\alpha\propto\Lambda_m$ so $\alpha_X=\alpha_Y/100$ and $K_a=C\alpha^2$ differs by $10^4$: $\mathrm{p}K_a(HX)-\mathrm{p}K_a(HY)=4$}, difficulty=4, diff_parts={K1 S2 I1 C0 F1}, diff_note={concepts: molar conductivity and ionisation | steps: degree of ionisation proportional to molar conductivity; ratio of Ka; difference of pKa}]
The molar conductivity of a solution of a weak acid HX (0.1 M) is 100 times smaller than the molar conductivity of a solution of a weak acid HY (0.1 M). If $\lambda^0_{X^-}=\lambda^0_{Y^-}$, the difference in their p$K_a$ values, p$K_a$(HX) $-$ p$K_a$(HY), is (consider degree of ionization of both acids to be $\ll1$).
\end{bankq}
\begin{bankq}[id=TOAS18-P2-CHEM-S1-Q04, source={TOAS - Test 18 (30 Oct 2022)}, paper=2, qno=4, type=single, subtopic={Cells}, answer={D}, answer_src=key, difficulty=2, diff_parts={K1 S1 I1 C0 F0}, diff_note={concepts: concentration cell standard potentials | steps: electrode potentials of anode and cathode; standard emf}]
For the given cell at $25^\circ$C, select the correct statement -
\par\noindent Given : $E^0_{\mathrm{Cl_2/Cl^-}}=1.38$ V
\bfig{224.4}{TOAS18_p2_chem_s1_q04_1.png}
\begin{qoptions}
\opt{A}{Electrode oxidation potential is negative for anodic compartment.}
\opt{B}{Standard oxidation potential is negative for anodic compartment.}
\opt{C}{Standard Emf of the cell is zero}
\opt{D}{All of these}
\end{qoptions}
\end{bankq}

\section*{TOPS papers}
\begin{bankq}[id=TOPS04-P2-CHEM-S1-Q15, source={TOPS - Test 04 (18 Dec 2022)}, paper=2, qno=15, type=single, subtopic={Cells}, answer={A}, answer_src=key, difficulty=4, diff_parts={K1 S2 I1 C1 F0}, diff_note={concepts: ionic mobility | steps: speed and potential gradient; effect of charge and hydration; time}]
Ionic mobility of an ion in an aqueous solution is given as

$\text{Ionic Mobility}=\dfrac{\text{Speed of an ion}}{\text{Potential gradient}}$

Assuming ions such as $\mathrm{Na^+}$, $\mathrm{Al^{3+}}$, $\mathrm{K^+}$ and $\mathrm{Ti^{4+}}$ are moving in aqueous solutions through same potential gradient the time required to reach the respective electrode from same starting point will be maximum for
\begin{qoptions}
\opt{A}{$\mathrm{Ti^{4+}}$}
\opt{B}{$\mathrm{K^+}$}
\opt{C}{$\mathrm{Na^+}$}
\opt{D}{$\mathrm{Al^{3+}}$}
\end{qoptions}
\end{bankq}

\section*{Star Review (Other) papers}
\begin{bankq}[id=STAR03-P1-CHEM-S1-Q01, source={Star Review (Other) - Test 03 (04 Dec 2023)}, paper=1, qno=1, type=single, subtopic={Nernst equation}, answer={C}, answer_src=key, difficulty=4, diff_parts={K1 S2 I1 C1 F0}, diff_note={concepts: emf of concentration cells | steps: Nernst equation; effect of each concentration; compare}]
The emf of following three galvanic cells are found to be as $E_1$, $E_2$ and $E_3$ volts.\par
(i) $\mathrm{Zn\,|\,Zn^{2+}\,(0.01M)\,||\,Cu^{2+}\,(1M)\,|\,Cu}$\par
(ii) $\mathrm{Zn\,|\,Zn^{2+}\,(1M)\,||\,Cu^{2+}(0.01\ M)\,|\,Cu}$\par
(iii) $\mathrm{Zn\,|\,Zn^{2+}(1M)\,||\,Cu^{2+}\,(2M)\,|\,Cu}$\par
Choose the correct option :
\begin{qoptions}
\opt{A}{$E_1>E_2>E_3$}
\opt{B}{$E_3>E_2>E_1$}
\opt{C}{$E_1>E_3>E_2$}
\opt{D}{$E_2>E_3>E_1$}
\end{qoptions}
\end{bankq}
\begin{bankq}[id=STAR03-P1-CHEM-S2-Q05, source={Star Review (Other) - Test 03 (04 Dec 2023)}, paper=1, qno=5, type=multi, subtopic={Conductance}, answer={B,C,D}, answer_src=key, difficulty=4, diff_parts={K1 S2 I1 C0 F1}, diff_note={concepts: conductivity and its dependence | steps: intensive and extensive; limiting conductivity; metallic conduction; dilution effect}]
Which of the following statements is/are correct ?
\begin{qoptions}
\opt{A}{Specific conductance is an intensive property while molar conductivity is an extensive property of solution.}
\opt{B}{The limiting equivalent conductivity of weak electrolyte can not be determined exactly by extrapolation of the plot of $\Lambda_{eq}$ against $\sqrt C$ [Given : $C$ = concentration of weak electrolyte]}
\opt{C}{The conductance of metals is due to the movements of free electrons.}
\opt{D}{Molar conductance of $0.01$ KCl aqueous solution is more than that of $0.1$ M KCl aqueous solution}
\end{qoptions}
\end{bankq}
\begin{bankq}[id=STAR03-P2-CHEM-S4-Q09, source={Star Review (Other) - Test 03 (04 Dec 2023)}, paper=2, qno=9, type=numerical, subtopic={Electrolysis}, answer={127}, answer_src=key, difficulty=6, diff_parts={K1 S2 I1 C2 F1}, diff_note={concepts: volume of gases from electrolysis of dilute sulphuric acid | steps: charge passed; moles of electrons; moles of each gas; total volume at standard conditions}]
A dilute solution of sulphuric acid is electrolysed using a current of $0.10$ A for $2$ hours to produce hydrogen and oxygen gas. The total volume of gases produced at STP is \underline{\hspace{1cm}} $\mathrm{cm^3}$. (Nearest integer)\par
[Given : Faraday constant $F=96500\ \mathrm{C\,mol^{-1}}$ at STP, molar volume of an ideal gas is $22.7\ \mathrm{L\,mol^{-1}}$]
\end{bankq}

\section*{Nurture (General) papers}
\begin{bankq}[id=NUR01-P1-CHEM-S1-Q05, source={Nurture (General) - Test 01 (12 Jun 2021)}, paper=1, qno=5, type=multi, subtopic={Nernst equation}, answer={A,B,D}, answer_src=solved, working={E=0.18+0.06x6=0.54 V; hydrogen is oxidised at the anode; E(Ag+)=0.18+0.6=0.78 V}, difficulty=7, diff_parts={K2 S2 I2 C1 F1}, diff_note={concepts: hydrogen silver chloride cell, nernst equation | steps: cell reaction; Nernst equation; electrode roles; Ksp relation}]
Select the correct statement(s) regarding following galvanic cell taken at temperature $T$K.\par
$\mathrm{Pt\,|\,H_2(g,\,1\,bar)\,|\,HCl\,(10^{-3}\,M)\,|\,AgCl(s)\,|\,Ag(s)}$\par
Given that : $\dfrac{2.303RT}{F}=0.06$ volt\par
\& $E^\circ_{\mathrm{Ag|AgCl|Cl^-}}=-0.18$ volt
\begin{qoptions}
\opt{A}{Cell potential of this cell is 0.54 volt.}
\opt{B}{During discharging, $\mathrm{H_2(g)}$ converts into $\mathrm{H^+}$ (aq.)}
\opt{C}{$\mathrm{H_2(g)}$ is obtained at cathode during working of this galvanic cell.}
\opt{D}{$E^\circ_{\mathrm{Ag^+|Ag}}$ is equal to 0.78 volt if $K_{sp}$ of AgCl is $1\times10^{-10}$.}
\end{qoptions}
\end{bankq}
\begin{bankq}[id=NUR03-P1-CHEM-S1-Q01, source={Nurture (General) - Test 03 (27 Sep 2023)}, paper=1, qno=1, type=single, subtopic={Cells}, answer={A}, answer_src=solved, working={DeltaG=-130.79+109.56=-21.23 kJ so E0=0.22 V; E=0.22+0.0591x4=0.456 V}, difficulty=5, diff_parts={K2 S2 I1 C1 F0}, diff_note={concepts: emf of a hydrogen silver chloride cell, cells | steps: free energy change; standard emf; Nernst equation}]
Calculate the cell EMF in mV for\par
Pt $|$H$_2$(1 atm)$|$HC$\ell$(0.01 M)$|$AgC$\ell$(s)$|$Ag(s) at 298K\par
If $\Delta G^\circ_f$ values are at 25$^\circ$C.\par
$-109.56\,\dfrac{\text{kJ}}{\text{mol}}$ for AgC$\ell$(s) and $-130.79\,\dfrac{\text{kJ}}{\text{mol}}$ for (H$^+$ + C$\ell^-$)(aq)
\begin{qoptions}
\opt{A}{456 mV}
\opt{B}{654 mV}
\opt{C}{546 mV}
\opt{D}{None of these}
\end{qoptions}
\end{bankq}
\begin{bankq}[id=NUR03-P1-CHEM-S1-Q02, source={Nurture (General) - Test 03 (27 Sep 2023)}, paper=1, qno=2, type=single, subtopic={Nernst equation}, answer={C}, answer_src=solved, working={Zero emf when Ksp(AgBr)/[Br]=Ksp(AgCl)/[Cl], so [Br]/[Cl]=1:200}, difficulty=6, diff_parts={K2 S2 I2 C1 F0}, diff_note={concepts: concentration cell from sparingly soluble salts, nernst equation | steps: electrode potentials with Ksp; equate; ratio}]
Consider the cell Ag(s) $|$ AgBr(s) $|$ Br$^-$(aq) $||$ AgC$\ell$(s) $|$ C$\ell^-$(aq)$|$ Ag(s) at 25$^\circ$C. The solubility product constants of AgBr \& AgC$\ell$ are respectively $5\times10^{-13}$ \& $1\times10^{-10}$. For what ratio of the concentrations of Br$^-$ \& C$\ell^-$ ions would emf of the cell be zero?
\begin{qoptions}
\opt{A}{1:500}
\opt{B}{1:100}
\opt{C}{1:200}
\opt{D}{200:1}
\end{qoptions}
\end{bankq}
\begin{bankq}[id=NUR03-P1-CHEM-S1-Q03, source={Nurture (General) - Test 03 (27 Sep 2023)}, paper=1, qno=3, type=single, subtopic={Cells}, answer={A}, answer_src=solved, working={E0=(0.77-2x0.44)/3=-0.0367 V}, difficulty=3, diff_parts={K1 S2 I1 C0 F0}, diff_note={concepts: potential of a combined couple | steps: free energies of the steps; combine; divide by electrons}]
If $E^\circ_{\mathrm{Fe}^{+2}/\mathrm{Fe}}=-0.44\,$V and $E^\circ_{\mathrm{Fe}^{+3}/\mathrm{Fe}^{+2}}=0.77$. What is $E^\circ_{\mathrm{Fe}^{+3}/\mathrm{Fe}}$
\begin{qoptions}
\opt{A}{-0.036V}
\opt{B}{-0.33 V}
\opt{C}{+0.33V}
\opt{D}{+0.0367V}
\end{qoptions}
\end{bankq}
\begin{bankq}[id=NUR03-P1-CHEM-S1-Q04, source={Nurture (General) - Test 03 (27 Sep 2023)}, paper=1, qno=4, type=single, subtopic={Electrolysis}, answer={D}, answer_src=solved, working={2 F per mole at 75 percent: i=2x96500/(0.75x3600)=71.48 A}, difficulty=3, diff_parts={K1 S2 I1 C0 F0}, diff_note={concepts: current for persulphate production | steps: electrons per mole; efficiency; current}]
Electrolysis of dil. H$_2$SO$_4$ generates H$_2$S$_2$O$_8$ (per disulphuric acid). What current is needed to produce H$_2$S$_2$O$_8$ at the rate of 1 mole per hour assuming 75\% charge efficiency.
\begin{qoptions}
\opt{A}{53.7amp}
\opt{B}{17.48 amp}
\opt{C}{35.74 amp}
\opt{D}{71.48 amp}
\end{qoptions}
\end{bankq}
\begin{bankq}[id=NUR03-P1-CHEM-S1-Q05, source={Nurture (General) - Test 03 (27 Sep 2023)}, paper=1, qno=5, type=single, subtopic={Nernst equation}, answer={C}, answer_src=solved, working={pH differs by (5-3)/2=1 so emf=0.059 V}, difficulty=6, diff_parts={K2 S2 I2 C1 F0}, diff_note={concepts: cell of two weak acid hydrogen electrodes, nernst equation | steps: hydrogen ion concentrations; Nernst equation; emf}]
Two weak acid solution HA$_1$ and HA$_2$ each with the same concentration and having pK$_a$ values 3 and 5 are placed in contact with hydrogen electrodes (1atm,25$^\circ$C) and are interconnected through a salt bridge. The emf of the cell is: (Assuming $\alpha\ll1$ for both acids)\par
\bfig{103.5}{NUR03_p1_chem_s1_q05_1.png}
\begin{qoptions}
\opt{A}{0.21V}
\opt{B}{0.018V}
\opt{C}{0.059V}
\opt{D}{0.021 V}
\end{qoptions}
\end{bankq}
\begin{bankq}[id=NUR03-P1-CHEM-S1-Q06, source={Nurture (General) - Test 03 (27 Sep 2023)}, paper=1, qno=6, type=single, subtopic={Electrolysis}, answer={D}, answer_src=solved, working={H2 0.434 mol gives 0.868 mol e-; O2 0.105 mol uses 0.42 mol e-; 0.448 mol e- gives 0.224 mol of H2S2O8 = 43.47 g}, difficulty=6, diff_parts={K2 S2 I2 C1 F0}, diff_note={concepts: persulphate and oxygen at the electrodes, electrolysis | steps: electrons for hydrogen; electrons for oxygen; remaining for persulphate; mass}]
During the preparation of H$_2$S$_2$O$_8$ (per disulphuric acid) O$_2$ gas also releases at anode as by product, When 9.72 L of H$_2$ releases at cathode and 2.35 L O$_2$ at anode at STP, the weight of H$_2$S$_2$O$_8$ produced in gram is
\begin{qoptions}
\opt{A}{87.12}
\opt{B}{51.74}
\opt{C}{83.42}
\opt{D}{43.47}
\end{qoptions}
\end{bankq}
\begin{bankq}[id=NUR03-P1-CHEM-S1-Q07, source={Nurture (General) - Test 03 (27 Sep 2023)}, paper=1, qno=7, type=single, subtopic={Nernst equation}, answer={B}, answer_src=solved, working={The weakest acid gives the lowest H+ concentration: acetic acid}, difficulty=2, diff_parts={K1 S1 I1 C0 F0}, diff_note={concepts: hydrogen electrode in different acids | steps: hydrogen ion concentration of each; highest emf}]
A galvanic cell is composed of two hydrogen electrodes, one of which is a standard one. In which of the following solutions should the other electrode be immersed to get maximum emf ? K$_a$(CH$_3$COOH) $=2\times10^{-5}$, K$_a$(H$_3$PO$_4$) $=10^{-3}$
\begin{qoptions}
\opt{A}{0.1 M HCl}
\opt{B}{0.1 M CH$_3$COOH}
\opt{C}{0.1 M H$_3$PO$_4$}
\opt{D}{0.1 M H$_2$SO$_4$}
\end{qoptions}
\end{bankq}
\begin{bankq}[id=NUR03-P1-CHEM-S1-Q08, source={Nurture (General) - Test 03 (27 Sep 2023)}, paper=1, qno=8, type=single, subtopic={Conductance}, answer={A}, answer_src=solved, working={Strong acid with strong base: conductance falls then rises, V shaped with the minimum at the end point}, difficulty=2, diff_parts={K1 S1 I1 C0 F0}, diff_note={concepts: conductance titration curve | steps: ions that change; shape of the curve}]
Conductance measurements can be used to detect the end point of acid-base titrations. Which of the following plots correctly represent the end point of the titration of strong acid and a strong base ?
\begin{qoptions}
\opt{A}{\bfigi{85.9}{NUR03_p1_chem_s1_q08_A.png}}
\opt{B}{\bfigi{84.1}{NUR03_p1_chem_s1_q08_B.png}}
\opt{C}{\bfigi{86.9}{NUR03_p1_chem_s1_q08_C.png}}
\opt{D}{\bfigi{81.9}{NUR03_p1_chem_s1_q08_D.png}}
\end{qoptions}
\end{bankq}
\begin{bankq}[id=NUR03-P1-CHEM-S1-Q09, source={Nurture (General) - Test 03 (27 Sep 2023)}, paper=1, qno=9, type=single, subtopic={Cells}, answer={A}, answer_src=solved, working={Pb and PbO2 both turn into PbSO4 on discharge}, difficulty=1, diff_parts={K1 S1 I0 C0 F0}, diff_note={concepts: lead accumulator discharge | steps: reactions at the plates}]
When a lead storage battery is discharged?
\begin{qoptions}
\opt{A}{PbSO$_4$ is formed}
\opt{B}{Pb is formed}
\opt{C}{SO$_2$ is consumed}
\opt{D}{H$_2$SO$_4$ is formed}
\end{qoptions}
\end{bankq}
\begin{bankq}[id=NUR03-P1-CHEM-S1-Q10, source={Nurture (General) - Test 03 (27 Sep 2023)}, paper=1, qno=10, type=single, subtopic={Cells}, answer={D}, answer_src=solved, working={PbO2 (1.455 V) cannot oxidise Mn2+ to permanganate (1.51 V)}, difficulty=2, diff_parts={K1 S1 I1 C0 F0}, diff_note={concepts: oxidation by couples of different potential | steps: compare the potentials; find the wrong statement}]
The standard reduction potentials of some half cell reactions are given below :\par
PbO$_2$ + 4H$^+$ + 2e$^-\rightleftharpoons$ Pb$^{2+}$ + 2H$_2$O \quad $E^0=1.455$ V\par
MnO$_4^-$ + 8H$^+$ + 5e$^-\rightleftharpoons$ Mn$^{2+}$ + 4H$_2$O \quad $E^0=1.51$ V\par
Ce$^{4+}$ + e$^-\rightleftharpoons$ Ce$^{3+}$ \quad $E^0=1.61$ V\par
H$_2$O$_2$ + 2H$^+$ + 2e$^-\rightleftharpoons$ 2H$_2$O \quad $E^0=1.71$ V\par
Pick out the \textbf{incorrect} statement (for standard condition of all species):
\begin{qoptions}
\opt{A}{Ce$^{+4}$ will oxidise Pb$^{2+}$ to PbO$_2$}
\opt{B}{MnO$_4^-$ will oxidise Pb$^{2+}$ to PbO$_2$}
\opt{C}{H$_2$O$_2$ will oxidise Mn$^{+2}$ to MnO$_4^-$}
\opt{D}{PbO$_2$ will oxidise Mn$^{+2}$ to MnO$_4^-$}
\end{qoptions}
\end{bankq}
\begin{bankq}[id=NUR03-P1-CHEM-S1-Q11, source={Nurture (General) - Test 03 (27 Sep 2023)}, paper=1, qno=11, type=single, subtopic={Electrolysis}, answer={C}, answer_src=solved, working={A copper anode dissolves instead of liberating oxygen}, difficulty=2, diff_parts={K1 S1 I1 C0 F0}, diff_note={concepts: oxygen at the anode | steps: reactive electrodes; discharge of ions}]
On electrolysis, in which of the following, O$_2$ would \textbf{not} be liberated at the anode ?
\begin{qoptions}
\opt{A}{dilute H$_2$SO$_4$ with Pt electrodes}
\opt{B}{aqueous AgNO$_3$ solution with Pt electrodes}
\opt{C}{dilute H$_2$SO$_4$ with Cu electrodes}
\opt{D}{aqueous NaOH with a Fe cathode \& a Pt anode}
\end{qoptions}
\end{bankq}
\begin{bankq}[id=NUR03-P1-CHEM-S1-Q17, source={Nurture (General) - Test 03 (27 Sep 2023)}, paper=1, qno=17, type=single, subtopic={Conductance}, answer={C}, answer_src=solved, working={Lambda0=138.2; S=1000x1.382e-6/138.2=1e-5 M}, difficulty=4, diff_parts={K2 S2 I1 C0 F0}, diff_note={concepts: solubility from conductivity, conductance | steps: limiting molar conductivity; relation; solve}]
The conductivity of a solution of AgCl at 298 K is found to be $1.382\times10^{-6}\,\Omega^{-1}$cm$^{-1}$. The ionic conductance of Ag$^+$ and Cl$^-$ at infinite dilution are $61.9\ \Omega^{-1}$cm$^2$mol$^{-1}$ and $76.3\ \Omega^{-1}$cm$^2$mol$^{-1}$, respectively. The solubility of AgCl is
\begin{qoptions}
\opt{A}{$1.4\times10^{-6}$mol L$^{-1}$}
\opt{B}{$1\times10^{-2}$mol L$^{-1}$}
\opt{C}{$1\times10^{-5}$mol L$^{-1}$}
\opt{D}{$1.9\times10^{-5}$mol L$^{-1}$}
\end{qoptions}
\end{bankq}
\begin{bankq}[id=NUR03-P1-CHEM-S1-Q18, source={Nurture (General) - Test 03 (27 Sep 2023)}, paper=1, qno=18, type=single, subtopic={Conductance}, answer={D}, answer_src=solved, working={BaSO4 = BaCl2 + H2SO4 - 2 HCl; equivalent conductance is half}, difficulty=2, diff_parts={K1 S1 I1 C0 F0}, diff_note={concepts: equivalent conductance by Kohlrausch | steps: combination of salts; equivalent factor}]
Molar conductances of BaCl$_2$, H$_2$SO$_4$ and HCl at infinite dilutions are x$_1$, x$_2$ and x$_3$ respectively. Equivalent conductance of BaSO$_4$ at infinite dilution will be:
\begin{qoptions}
\opt{A}{$\dfrac{[x_1+x_2-x_3]}{2}$}
\opt{B}{$\dfrac{[x_1-x_2-x_3]}{2}$}
\opt{C}{$2\,(x_1+x_2-2x_3)$}
\opt{D}{$\dfrac{[x_1+x_2-2x_3]}{2}$}
\end{qoptions}
\end{bankq}
\begin{bankq}[id=NUR03-P1-CHEM-S1-Q19, source={Nurture (General) - Test 03 (27 Sep 2023)}, paper=1, qno=19, type=single, subtopic={Electrolysis}, answer={D}, answer_src=solved, working={Without a complete circuit no field acts and the ions move randomly}, difficulty=2, diff_parts={K1 S1 I1 C0 F0}, diff_note={concepts: ions without a cathode | steps: electric field on the ions}]
When the electric current is passed through a cell having an electrolyte, the positive ions move towards cathode and negative ions towards the anode. If the cathode is pulled out of the solution
\begin{qoptions}
\opt{A}{The positive ions will move towards anode}
\opt{B}{The positive ions will start moving towards the anode while negative ions will stop moving}
\opt{C}{The negative ions will continue to move towards anode while positive ions will stop moving}
\opt{D}{The positive and negative ions will start moving randomly}
\end{qoptions}
\end{bankq}
\begin{bankq}[id=NUR03-P1-CHEM-S2-Q03, source={Nurture (General) - Test 03 (27 Sep 2023)}, paper=1, qno=3, type=numerical, subtopic={Nernst equation}, answer={0.32}, answer_src=solved, working={E0 of the complex/Ag = 0.80-0.06x8/... =0.32 V reduction; magnitude of the oxidation potential 0.32 V}, difficulty=7, diff_parts={K2 S2 I2 C1 F1}, diff_note={concepts: oxidation potential of a complex couple, nernst equation | steps: formation constant; Nernst equation; reduction potential}]
The magnitude of $E^\circ_{\mathrm{OP}}$ of following reaction is \underline{\hspace{1cm}}volt.\par
Ag(s) + 2NH$_3$ (aq) $\to$ Ag(NH$_3$)$_2^+$ (aq) + e$^-$\par
[Given : K$_f$ (Ag(NH$_3$)$_2^+$) $=10^8$; $E^\circ_{\mathrm{Ag}^+|\mathrm{Ag}}=0.80$V; $2.303\dfrac{RT}{F}=0.06$]
\end{bankq}
\begin{bankq}[id=NUR03-P1-CHEM-S2-Q04, source={Nurture (General) - Test 03 (27 Sep 2023)}, paper=1, qno=4, type=numerical, subtopic={Cells}, answer={10}, answer_src=solved, working={dE/dT=4e-4 V/K, E(300)=0.068 V; deltaH=-nFE+nFT dE/dT with n=2: 10 kJ}, difficulty=8, diff_parts={K3 S2 I2 C1 F1}, diff_note={concepts: enthalpy from the emf temperature graph, cells, slope of emf | steps: slope of emf; emf at 300 K; Gibbs Helmholtz relation}]
For the cell reaction: Hg$_2$Br$_2$ (s) + 2Ag (s) $\to$ 2Hg ($\ell$) + 2AgBr (s) ,the variation of emf with temperature is shown by the graph given below\par
\bfig{141.4}{NUR03_p1_chem_s2_q04_1.png}
Calculate $\Delta H$ of cell reaction(in kJ) at 300K.
\end{bankq}
\begin{bankq}[id=NUR03-P1-CHEM-S2-Q05, source={Nurture (General) - Test 03 (27 Sep 2023)}, paper=1, qno=5, type=numerical, subtopic={Nernst equation}, answer={1.48}, answer_src=solved, working={0.79-0.12pH=0.731-0.08pH gives pH=1.475}, difficulty=7, diff_parts={K2 S2 I2 C1 F1}, diff_note={concepts: pH where two nitrate couples have equal potential, nernst equation | steps: half reactions; Nernst equation for each; equate}]
NO$_3^-\to$ NO$_2$(acid medium), E$^\text{o}$=0.790V\par
NO$_3^-\to$ NH$_3$OH$^+$(acid medium) E$^\text{o}$=0.731V.\par
At what pH. the above two will have same E value? Assume the concentration of all other species NH$_3$OH$^+$ except[H$^+$] to be unity.
\end{bankq}
\begin{bankq}[id=NUR03-P1-CHEM-S2-Q06, source={Nurture (General) - Test 03 (27 Sep 2023)}, paper=1, qno=6, type=numerical, subtopic={Electrolysis}, answer={4}, answer_src=solved, working={0.02798 mol of Pd needs 10800 C = 386000 C per mole = 4 F}, difficulty=4, diff_parts={K1 S2 I1 C0 F1}, diff_note={concepts: valency from electrolysis | steps: charge passed; moles deposited; electrons per ion}]
3A current was passed through an aqueous solution of an unknown salt of Pd$^{n+}$ for 1Hr. 2.977g of Pd was deposited at cathode. Find n.(Given Atomic mass of Pd=106.4)
\end{bankq}
\begin{bankq}[id=NUR03-P1-CHEM-S2-Q07, source={Nurture (General) - Test 03 (27 Sep 2023)}, paper=1, qno=7, type=numerical, subtopic={Cells}, answer={1}, answer_src=solved, working={Reduction potentials: A -0.92, B -0.18, C +0.02, D +1.20; C as anode only with D as cathode}, difficulty=7, diff_parts={K2 S2 I2 C1 F1}, diff_note={concepts: feasible cells with a given anode, convert to reduction potentials | steps: convert to reduction potentials; compare; count}]
Based on the following standard oxidation electrode\par
$E^\circ_{\mathrm{A}^{+2}/\mathrm{A}}=-0.92$V\par
$E^\circ_{\mathrm{B}/\mathrm{B}^{+1}}=+0.18$V\par
$E^\circ_{\mathrm{C}^{+3}/\mathrm{C}}=+0.02$V\par
$E^\circ_{\mathrm{D}/\mathrm{D}^{+2}}=-1.20$V\par
Number of feasible electrochemical cells that can be constructed using metal C as anode at 298K will be
\end{bankq}
\begin{bankq}[id=NUR03-P1-CHEM-S2-Q08, source={Nurture (General) - Test 03 (27 Sep 2023)}, paper=1, qno=8, type=numerical, subtopic={Cells}, answer={45}, answer_src=solved, working={E0(Fe3+/Fe2+)=3(-0.036)+0.88=0.772 V; E0cell=0.233 V; deltaG=-2x96500x0.233=-45 kJ}, difficulty=7, diff_parts={K2 S2 I2 C1 F1}, diff_note={concepts: free energy of the iron iodide reaction, cells | steps: potential of Fe3+/Fe2+; cell potential; free energy}]
For the reaction 2Fe$^{3+}$(aq) + 2I$^-$(aq) $\to$ 2Fe$^{2+}$(aq) + I$_2$(s)\par
The magnitude of the standard molar free energy change, $\Delta_rG^\circ_m=-$\underline{\hspace{1.5cm}} kJ (Round off to the Nearest Integer).\par
[$E^\circ_{\mathrm{Fe}^{2+}/\mathrm{Fe(s)}}=-0.440$ V; $E^\circ_{\mathrm{Fe}^{3+}/\mathrm{Fe(s)}}=-0.036$ V; $E^\circ_{\mathrm{I}_2/2\mathrm{I}^-}=0.539$ V; $F=96500$ C]
\end{bankq}
\end{document}
